\section{Field lists}
\label{supp:fields}

This section lists, entry by entry, the fields of the records that the main paper names: the
certificates of the Templates, the status branches of each conditional classification, and the data
of the definitions. A field is a hypothesis of the corresponding Lean statement; nothing listed here
is derived from the type of an E-system. Throughout, a \emph{claim} is the full tuple that a status
record is about (system, challenge class, horizon, and the records the entry names), and every field
is scoped to that claim: evidence for one claim cannot be used for another.

\subsection{The definitional layer}

\paragraph{Carried record (D3).} A declared coordinate map $\rho_{\mathrm{of}} : Z \to R$; a
declared trajectory $\tau$ with a legitimate start; an index $n_0 \ge n_{\mathrm{start}}$ such that
every step from $n_{\mathrm{start}}$ to $n_0$ is in scope and lies in $\supp K$; the value
$\rho = \rho_{\mathrm{of}}(\tau(n_0))$; a source tag among committed state and audited cell records;
generation and scope flags equal to true. The source tags are committed state, audited cell records,
independent-pair witness, simulation trace, ablation record, fallback, unknown and contradictory.

\paragraph{Carried instrument (D3).} An instrument of the audit calculus of \FoundIII{}; a predicate
saying which instrument records are carried; lists of visibility, threshold and check-rule records,
each carried; and a proposition declaring the lists complete.

\paragraph{E-system (D4).} A theory package with carrier, kernel support, legitimate-start predicate,
declared trajectory, scope predicate, lens, definability structure, completion, audit, and a
formedness proposition with its proof; carried-record policies for defect, move and audit records; an
active carried instrument with the predicates \emph{detects}, \emph{gate allows} and
\emph{re-audits}; a carried ledger (a policy, a list of entries, a completeness proposition with its
proof, carriedness of every entry); a repair generator sending each defect to a move with a sort in
P1--P5, a payload, a carried move record and a budget line; and an admissibility predicate.

\paragraph{Lawful repair step (D4).} Defect carried; detected at $z$; gate allows the generated move
at $z$; move admissible at $z$; budget line in the ledger; $(z, z') \in \supp K$; re-audit at $z'$
with the audit record; audit record carried.

\paragraph{Challenge history and closed-loop scope (D5).} Episodes, each with a time, a challenge
class, a source quotient, a declared family, a target readout and an optional obstruction witness;
repair-typed entries (defect, move or audit), each with an entry class carrying an identifier, a
source tag and generation and scope flags; an association of entries to episodes, declared for every
listed entry; a completeness proposition. Closed-loop scope is completeness together with
carriedness of every listed entry.

\paragraph{Probe economy (D6).} A finite catalog; a carried active family with a support list, a
weight function and an optional operator family; a carried exposure budget and spend in the ledger;
a saturation predicate; allocation, acquisition and retirement moves, each with pre- and post-family
carriedness records, and for acquisition strictness, catalog membership, budget admissibility and
absence from the prior support.

\subsection{Endogeny}

\paragraph{E1.} The endogenous branch contains, for each challenged episode: the defect record and
its generated move; pre- and post-states; the audit record; the repair payload and the installed
refinement with $Q_{t+1} = Q_t \vee R_t$; a strict challenged-fibre reduction of the split-pair
obstruction; a committed-state source; kernel realization; gate passage; and re-audit. The other
branches are: slack (the challenge is not binding); externally subsidized (a discharging entry
that is not endogenous and no valid endogenous family); stressed (viability preserved while a
statused residual accrues); collapsing (a viability probe fails). The certified inputs are
no-free-distinction, finite-forcing strictness, promotion soundness, the bridge from measured
discharge to absence of residual accrual, and family-level budget feasibility.

\paragraph{E2 (first version).} A carried instrument level has a level index, a tag, an instrument,
its own instrument record, an audit predicate on lists of lower indices, a lower-stack target, and the
ledger entries and audit records it uses. It occurs if its own record is carried by its instrument and
by the system instrument, its instrument is carried, its ledger entries are in the ledger, its audit
records are carried, and every target index is below its own. The first-version measure has injective,
disjoint tags for the three record sorts; a record universe per state with size equal to the capacity;
for every occurring level a record set of exactly its counted records, inside every universe, with
positive size; and disjointness between occurring levels of different index.

\paragraph{E2 (declared measure).} As in \Cref{def:declared-measure} of the main paper: depth, a
duplicate-free slot list covering every index below the depth, a level at each slot with matching
index and a proof that the level occurs, the tag data, the universes and capacities, the slot record sets with
the counting, inclusion, disjointness and positivity fields. A declared tower lists the level of each
slot at its index.

\paragraph{E2 (forced stratification).} An audited-repair trigger; the audit-stack record of
\FoundIII{} with finite records and an admissible lower stack; a bridge from the tower to that stack;
non-saturation; and a fresh-level bridge requiring every carried level at the target index to carry a
tag above every tag below it and to audit exactly the lower stack.

\paragraph{E3.} Apparatus records at $t$ and $t+1$ with their carried records (time, gate instrument,
threshold records, boundary, audit data, apparatus record, used ledger entries and audit records); a
maintenance operator with its record, its action on states, and its ledger entries and audit records;
a reinstatement record (time, source and target states, operator, pre-, post- and output apparatus
records, ledger entry) that is carried with an allowed source, whose fields match, whose post
apparatus is the operator applied to the source, whose ledger entry is in the ledger and whose source
and target are joined by a kernel move. The maintained branch adds an object fixed point, object
record coherence, a maintained step (distance at most $\varepsilon$ without challenge; a linked E1
repair with defect, move, join installation, kernel transition and re-audit under challenge), an empty
maintenance obstruction, no subsidy, and regress-stop evidence. The maintenance obstruction is
nonempty if a required component (instrument, gates, thresholds, boundary, audit data) is absent,
has a forbidden source, fails generation or scope, is off-kernel, or has a valid reinstatement omitted
from the audit.

\paragraph{E4.} A candidate window of at least two carried invocations of the same attributed
repair; an obstruction-reduction predicate and an idempotence-stability predicate, independent of
each other; a compiled operator record of sort P1 or P2, indexed by the challenge class and the repair
refinement; the lawfulness gates (certified top-down channel, accepted promotion record, failed
memory-only comparator, non-saturation, annotation record); verified in-class descent; silence of the
higher invocation. Statuses in priority order: decompiled (a linked decompilation event),
mis-compiled (the compiled core with an unstatused out-of-class obstruction), compiled, compiling
(an attributed candidate missing a gate, descent or silence), exposed (no attributed candidate).

\paragraph{E5.} Rescue candidates of four types (repair, boundary update, apparatus reclosure, probe
acquisition), each with a move record, a budget witness, a kernel transition and a descent
certificate for the declared viability probes; a list containing every candidate. Irreversibility
inputs: unreachable repair generator and empty apparatus-level viability kernel, coherently scoped.
Statuses in priority order: revived (a prior subsidized or collapsed status, a later carried rescue, a
new apparatus lineage), subsidized (an active external carrier with its own scope, budget and
descending move, and no internal falsifier), suspended (intact apparatus, operations gated off),
viable, stressed (additionally, scoped residual accrual), collapsed-irreversible and
collapsed-recoverable (the latter with its own recoverability evidence).

\subsection{Priced access and control}

\paragraph{E6.} A budget record: spend, budget, their ledger entries (in the ledger and certified by
the economy) and budget admissibility of the move. A multiplier witness: $\lambda \ge 0$, $\mu_p \ge
0$, the two slackness equations and stationarity on the support. A binding budget is spend $=$
budget with spend $\le$ budget; a slack budget is spend $<$ budget.

\paragraph{E7.} Currentization: a carried constraint rewrite linking old and new constraint sets, a
forced-exposure ledger entry naming the witness, the new set different from the old and admitting
the witness, and passing no-overread and null-mode checks. Boundary nonclosure: a failed viability
descent and a carried status. Lawful discount: a carried reason among protocol artifact,
currentizable residue, out of horizon, dissolved coupling. The package requires one branch-valid
carried record and the same kind for every carried record at the witness and horizon.

\paragraph{E8.} A certified signal readout; a declared constraint record; a ledger-alignment
comparator; held-out prediction data; a dual-stability record under declared budget perturbations;
a linked multiplier witness. A binding component: the constraint binds, the multiplier is the declared
dual price, alignment, held-out prediction and stability pass. A slack component: multiplier and
price component vanish, or a carried residual explains the deviation. Capture evidence: a linked budget
rise by the classified move, a positive blind spot against the original accepted lineage, and a healthy
ledger report; the audit statuses are capture-claim-rejected, capture-detected,
residual-discharge-verified, meta-audit-capacity-blocked and meta-audit-missing.

\paragraph{E9.} An admissible acquisition: catalog membership, absence from the support, strictness,
a budget record for the acquisition move with spend at most budget, and the viability-risk predicate.
A complete candidate list contains every admissible acquisition. An access move is an allocation
increment on the current support or a strict acquisition, priced by a common ratio.

\paragraph{E10.} Repair descent: a nonempty obstruction before, a quotient installed by the package,
and a strict reduction after. Gate: two distinct carried interventions, matched controls, a channel
record linked to that comparison, accepted gate status, and different induced kernel supports.
Provenance: carried formation and invocation records, the invocation linked to the class.
Statuses in priority order: schedule trap (an exogenous schedule explains the support change and a
complete gate inventory has no accepted gate), scaffolding (no carried package provenance),
structural path only (a lower-support difference and a structural path with the channel blocked),
decodable correlate, repair without gate, gate without repair, cognitive, non-cognitive. Intention:
provenance, a top-down restriction strictly before the outcome, strict inclusion of the restricted
action support, and coherence under every member of a nonempty carried list of future probes;
statuses post-hoc, approximate support, target-incoherent, intention, not intention. Goal: provenance,
top-down orientation, a carried target quotient, and evaluation of every listed route among at least
two distinct carried routes; statuses reward proxy only, route unstable, goal, not goal.

\subsection{Stack, individuation and transport}

\paragraph{E11.} Two carried intervention records and a matched-control record; a top-down channel
of \FoundIII{} referring to that comparison, accepted and passing its effect gate; a support witness
$(z, z')$ with different support indicators, or different closures. The stack obstruction is inhabited
by an independently carried parameter-only comparison (parameters differ, support and closure
unchanged) with the same outcome. Statuses in priority order: constitutive (structural effect, empty
stack obstruction, a carried compiled-rewrite record with a certified bridge), stack-active,
conditioning (parameter-only comparison and no structural effect), inert (washout, blocked channel,
or no evidence).

\paragraph{E12.} Repair closure: some repair is attributed to $I$ and every attributed repair has its
target in $I$. Record closure: some genuine activity is attributed to $I$ and every entry it generates
closes on $I$. Budget closure: a feasible exposure budget attributed through the carried trajectory,
every such entry closing on $I$. Boundary: a quotient certified sufficient for the declared viability
probes that is a function of every admitted sufficient quotient (so it is the coarsest), and
self-maintained by a boundary apparatus reinstated in the sense of E3 with source state in $I$.
Maximality ranges over the declared candidate list. Statuses: integrated; federated (a
non-individuating candidate properly inside a maximal one, or two overlapping maximal candidates);
subsidiary (repair and records close, budget does not); platform-dependent (all closures and a
sufficient coarsest boundary, not self-maintained); shadow (repair and budget close, records do not);
non-individuated.

\paragraph{E13.} A paid, carried bridge between the systems; a carried lawful receiver move; a strict,
installed reduction of the receiver's obstruction in the shared class; interface mediation; repair-role
preservation. Teaching adds a later probe with the source absent, the bridge closed and the repair
carried by the receiver from committed state. Coercion null: a comparator-certified fully forced
interaction, a carried, complete and empty free-response census, no receiver-committed transported repair, and
equal before and after obstructions. Symbolic repair: a paid carried bridge; a declared context family
with at least two entries; reactivation in every listed context; stability of the repair role;
same-family strictness. Statuses in priority order: coercion null, symbolic, taught, communicated,
scaffolded or primed, influence only, transport rejected.

\subsection{Memory and adaptive capacity}

\paragraph{E14.} A nonempty duplicate-free retrieval-context family; a complete transport inventory,
single-valued per context; context dependence (two contexts whose transports differ); the conflict, a
split pair after the transported record is joined to the current quotient; formed carried source
records. Repair: a lawful repair step, a carried next-version record, exact retrieval and context
linkage, honest provenance, installation of the record join. Coarsening: the before quotient refines the
after quotient, a strictness pair, and a carried audit that the lost
distinction is unsupported. Retained debt: unchanged record and a positive carried residual charged to
an existing ledger entry. Statuses in priority order: provenance defect, silent rewrite, ordinary read,
outcome collision, unrealized disposition, record repaired, record coarsened, statused unresolved,
unstatused conflict.

\paragraph{E15.} Closure debt from canonically linked E14 statused residuals, certified route residues
and adequacy residuals. The allocation record: carried, linked to the scope, declared before the
horizon, a binding budget with carried entries, a positive online external share, a nonnegative online
internal share, online exhaustion, offline exchange zero, the offline internal share equal to the
online internal plus external shares, offline exhaustion. Alternation required: strict online deficit,
recurrent lawful offline phases on the declared flow inventory with linked debt discharge, a
duty-cycle prediction, a tolerance declared before the horizon and the observed duty cycle within it.
Statuses in priority order: deficit-online counterexample, decorative offline, duty-cycle mismatch,
skipped-offline cascade, alternation required, online sufficient, offline optional, rejected.

\paragraph{E16.} A carried route-transport package with the comparison base, the transport core of the
holonomy development, two routes and their endpoints; current equivalence of the endpoints and a
predictive witness separating them; complete endpoint-conditioned trial inventories with eligibility,
success, protocol, measurement time and binding budget common to both arms and positive denominators;
a held-out challenge with $p_\gamma(C) \ne p_\eta(C)$ for the claim-linked pair; an accepted
reconciliation with any E15 route residue on the same package and pair; and four freedom certificates:
no eligible honest counterpart changes the comparison base; no admissible completion equalizes the
residue; no same-base refinement makes it current; no own continuation erases it and every declared
in-bound perturbation preserves it. Statuses in priority order: support confound, artifact, flat,
flattenable, currentizable slack, dissipative, coherent adaptability, adaptability rejected.
